\documentclass[10pt,a4paper]{amsart}
\usepackage[margin=19mm]{geometry}
\usepackage{amsmath,amssymb,mathtools}
\usepackage[scr=rsfs,cal=euler]{mathalfa}
\usepackage{xfrac}
\usepackage[unicode,hidelinks,bookmarks=false]{hyperref}
\newcommand{\ie}{i.e.\ }
\newcommand{\cf}{cf.\ }
\newcommand{\resp}{respectively\ }
\newcommand{\Subsection}{\subsection}
\providecommand{\nobreakdash}{\nobreak\hbox{-}}
\setlength{\emergencystretch}{2em}
\multlinegap0pt
\usepackage{bm}
\usepackage{enumitem}
\usepackage{centernot}
\usepackage{marginnote}
\usepackage{bbm}
\newtheorem{thm}{Theorem}[section]
\newtheorem{prop}[thm]{Proposition}
\newtheorem{lemma}[thm]{Lemma}
\newtheorem{cor}[thm]{Corollary}
\newtheorem{remark}[thm]{Remark}
\newtheorem{example}[thm]{Example}
\newtheorem{nota}[thm]{Notations}
\newtheorem{df}[thm]{Definition}
\newtheorem{q}[thm]{Question}
\newcommand{\R}{\mathbb{R}}
\newcommand{\norm}[1]{\left\Vert #1\right\Vert}
\begin{document}
\title[Stability of Individually Eventually Positive Semigroups on ]{Stability of Individually Eventually Positive Semigroups on $L^p$-Spaces}
\author{Loris Arnold}
\date{}
\maketitle
\noindent\emph{Source and licence.} Stability of Individually Eventually Positive Semigroups on $L^p$-Spaces. Version arXiv:2608.15038v1 (CC BY 4.0). This material is distributed under the Creative Commons Attribution 4.0 International licence, http://creativecommons.org/licenses/by/4.0/, as recorded in the arXiv OAI licence metadata for this exact version and shown on the abstract page https://arxiv.org/abs/2608.15038v1. Full rights notice: licenses/arxiv-2608.15038-CC-BY-4.0.txt.\par\smallskip
\noindent\textit{Editorial scope.} Counted unit: the original \texttt{thm} environment \texttt{thmRange} at source lines 190--211 together with its complete proof at lines 213--240; the delivered document keeps source lines 186--241 so that every referenced label and every citation resolves inside it. Native editable source is the paper's own concatenated TeX; the AMS wrapper is editorial formatting only. No publisher class, upstream script or unrelated artwork is redistributed.
	\section{A uniformisation principle for time-dependent operator ranges}
	
	We first present the operator-range principle that will be used below; an analogous statement can be found in \cite[Theorem 2.6]{AroraGlueck2024}.
	

	\begin{thm}
		\label{thmRange}
		Let $E$ and $F$ be Banach spaces, and let $(S_i)_{i\in I}$ be a net of bounded linear operators from $E$ to $F$. Assume that the directed set $I$ contains a countable majorizing subset. For each $i\in I$, let $W_i\subseteq F$ be a vector subspace endowed with a complete norm $\norm{\cdot}_{W_i}$ such that the inclusion $W_i\hookrightarrow F$ is continuous (that is, $W_i$ is an operator range in $F$). Then the following statements are equivalent.
		\begin{enumerate}[label=\upshape(\alph*)]
			\item\label{itUnif}
			There exist $M_0\ge 0$ and $i_0\in I$ such that, for every $i\ge i_0$,
			\[
			S_i E \subseteq W_i
			\qquad\text{and}\qquad
			\norm{S_i}_{B(E, W_i)} \le M_0.
			\]
			
			\item\label{itIndiv}
			There exists $M_1\ge 0$ such that for every $x\in E$ there exists $i_x\in I$ for which
			\[
			S_i x \in W_i
			\qquad\text{and}\qquad
			\norm{S_i x}_{W_i} \le M_1\norm{x}_E
			\]
			for every $i\ge i_x$.
		\end{enumerate}
	\end{thm}

	\begin{proof}
		The proof is essentially the same as that of \cite[Theorem 2.6]{AroraGlueck2024}. For the sake of completeness, we give a short proof. The implication \ref{itUnif}$\Rightarrow$\ref{itIndiv} is immediate.
		
		Conversely, let $I_0\subseteq I$ be a countable majorizing subset. For $i\in I_0$, set
		\[
		V_i := \left\{ x\in E \mid S_j x\in W_j \text{ for all } j\ge i \text{ and } \sup_{j\ge i}\norm{S_j x}_{W_j}<\infty \right\}
		\]
		and endow $V_i$ with the norm
		\[
		\norm{x}_{V_i} := \norm{x}_E + \sup_{j\ge i}\norm{S_j x}_{W_j}.
		\]
		
		A standard completeness argument shows that $V_i$ is an operator range in $E$. Indeed, if $(x_n)$ is Cauchy in $V_i$, then it converges in $E$, while for each $j\ge i$ the sequence $(S_j x_n)$ converges in $W_j$. The continuity of the embedding $W_j\hookrightarrow F$ and the boundedness of $S_j\colon E\to F$ identify this limit with $S_j x$.
		
		Assumption \ref{itIndiv} and the majorizing property of $I_0$ show that
		\[
		E = \bigcup_{i\in I_0} V_i.
		\]
		By \cite[Proposition 2.6]{AroraGlueck2023}, there exists $i_0\in I_0$ such that 
		\[
		V_{i_0} = E.
		\]
		The open mapping theorem then yields a constant $C\ge 0$ such that for every $x\in E$
		\[
		\sup_{j\ge i_0}\norm{S_j x}_{W_j} \le C\norm{x}_E.
		\]
		This is precisely \ref{itUnif}.
	\end{proof}
	

\begin{thebibliography}{99}
\bibitem[AroraGlueck2023]{AroraGlueck2023} S. Arora and J. Gl{\"u}ck, \emph{A characterization of the individual maximum and anti-maximum principle}, Math. Z. 305 (2023), art. 24.
\bibitem[AroraGlueck2024]{AroraGlueck2024} S. Arora and J. Gl\"uck, \emph{Irreducibility of eventually positive semigroups}, Stud. Math. \textbf{276} (2024), no. 2, 99--129.
\end{thebibliography}
\end{document}
